\section{State of the art}
\label{sec:state-of-the-art}

The literature relevant to the Research Map spans different units of work: formulating a question, generating a research idea, organizing a full project, and preserving the history of revisions. Similar interface features do not imply equivalent methodological purposes. This is a focused narrative review of seven supplied scholarly papers and the PROV-DM standard, complemented by public documentation for six applications; it is not a systematic search or a functional evaluation of those applications. We compare the artifact produced, the human researcher's role, the treatment of revisions, and the connection to the Research Map.

\subsection{From an initial problem to a research design}

Guerard et al. propose Research LLM as a three-stage human-in-the-loop framework covering proposal development, research design, and manuscript preparation \cite{guerard2026researchllm}. Their account combines prompting and verification techniques with explicit human checkpoints; it illustrates the pipeline through a hiring-bias case and reports an exploratory, qualitative pilot with master's students developing thesis proposals. The paper is close to the Research Map in treating the researcher as the decision maker across several stages. Its organizing artifact, however, is a broad research pipeline rather than a specified map linking a central problem, a general objective, specific objectives, chapter topics, and a collection--analysis--expected-result entry for each objective. The distinction concerns the structure being documented, not the presence of human review.

Wang et al. offer a more directly methodological comparator \cite{wang2026fivestage}. Their five-stage approach begins with observation and appraisal of a problem in general practice, transforms it into searchable concepts, summarizes the literature, constructs a research-question prototype, and then considers methodology and feasibility. Tables and prompts carry information between stages; clinicians inspect and amend intermediate outputs. The authors report feedback from a panel of ten general-practice researchers and one worked process log, rather than a controlled effectiveness trial. Both approaches connect a practical problem to literature and method, but their five stages are not the same five stages. Wang et al. organize the path toward a clinical research question; the Research Map organizes a project through problem, general and specific objectives, literature and development topics, with an accompanying objective--method matrix. The latter matrix asks separately what will be collected, how it will be analyzed, and what result is expected for each objective. An expected result is a planning entry, not an empirical finding.

CoQuest narrows the unit of assistance to research-question ideation \cite{liu2024coquest}. Its interface displays a flow of candidate questions, related literature, and the agent's rationale; users can explore breadth-first or depth-first alternatives and provide feedback. A within-subjects study with 20 HCI researchers compared those interaction designs and reported differences in perceived creativity and trust. CoQuest thus provides evidence about human--AI interaction during question generation that the present article cannot claim for the Research Map. Conversely, its evaluated task is not the maintenance of a project-wide objective--method plan. Research questions are relevant to the detailed Research Map method, but they are not a sixth numbered stage of its original five-stage structure; the current application also lacks a separate research-questions unit. This difference prevents us from treating CoQuest's user findings as evidence for our system.

\subsection{Agentic generation and critique}

ResearchAgent starts from a scientific paper and related literature, then proposes a problem, method, and experimental design; reviewing agents critique successive versions \cite{baek2025researchagent}. The authors compare model variants using human and model-based judgments of generated ideas. This work already joins literature grounding, methodological proposals, and iterative agent critique. Its primary artifact is an automatically generated research idea, whereas the Research Map is a researcher-maintained arrangement of linked project components. The comparison also exposes a shared risk: a plausible proposed method still requires substantive checking by a domain expert. ResearchAgent's evaluation results cannot be transferred to the Research Map, whose proposals, checkpoints, and outcomes have not been tested in that study.

Intelliscope, described by Aamer et al., connects biomedical entities through knowledge-graph paths and uses Analyst, Scientist, and Reviewer agents to explain those paths, design predictors and validation experiments, and refine proposals \cite{aamer2026biomedical}. Its ResearchMemory records design iterations and reviewer feedback, and its dashboard supports inspection of the generated work. It is a stronger comparator for retained agent decisions than a generic chatbot. Its knowledge-graph inputs and biomedical discovery outputs differ substantially from the Research Map's general project-planning fields. Consequently, neither multiple agent roles, retained feedback, nor human oversight can serve as the article's standalone novelty claim; the issue is how a particular methodological plan is represented and revised in the implemented workflow.

\subsection{Provenance and revision of structured artifacts}

VisTrails records actions that transform one scientific workflow into another and arranges workflow versions in a tree, allowing prior configurations and outputs to be revisited \cite{freire2006vistrails}. The W3C PROV-DM recommendation supplies a domain-independent account of entities, activities, agents, and derivations \cite{moreau2013prov}. RaivenTracks applies a more recent version-history design to conversational visualization: validated visualization specifications become branchable checkpoints, with a separate undo/redo layer for runtime settings \cite{hugie2026raiventracks}. RaivenTracks is a preprint; its formative pilot involved three visualization researchers and had no comparison baseline.

These sources clarify what the Research Map can mean by traceability. A record of a proposed field, its source revision, the researcher's decision, and the later version makes a change inspectable. It does not prove that the relation between a problem and an objective is scientifically sound. Unlike the workflow or visualization artifacts in VisTrails and RaivenTracks, the Research Map's objects are methodological statements and planning relations. In the current application, confirmed changes can mark downstream \emph{types} of component for review while preserving their text and history; this is conservative type-level invalidation, not an exact causal graph of every affected statement. PROV-DM is a conceptual comparator, not a claim that the application implements the W3C model in full.

\subsection{Applications documented outside the research literature}

The academic texts are not the whole landscape of practice. In this focused scan, ProjectReady AI is the closest publicly documented product comparator: its guide describes a persistent Research Record with approved objectives, questions, design decisions, researcher checkpoints, supervisor corrections, and an alignment audit \cite{projectready2026guide}. Researchly AI likewise advertises a path from research idea to plan, literature and method support, trust labels, supervisor review, and a ``what-if'' operation that recalculates parts of a plan after a change \cite{researchly2026site}. These are substantial overlaps with the Research Map's authoring and revision goals. Their public pages do not disclose enough internal behavior to determine whether their dependency rules are equivalent to the Research Map's, and we did not test authenticated workflows. It would therefore be unsound to claim either that these products lack such rules or that the Research Map is the first persistent AI-assisted research planner.

Other tools illustrate narrower or adjacent overlaps. Bodhi's Proposal Builder advertises generation of gaps, problem statements, objectives, hypotheses, methods, and expected outcomes from a topic \cite{bodhi2026proposal}. ResearchPlanBuilder's public repository describes guided planning questions, an optional AI step, a browser-local draft, and a generated research plan \cite{vallavaraj2026builder}. Reos plans user and product research through objectives, methods, participant criteria, questions, and a moderation guide, and its product page says that later evidence is linked back to those questions \cite{reos2026plan}. Tykyra structures a systematic literature review from protocol and search terms through screening and export, with claimed auditability of review decisions \cite{tykyra2026site}. The first two overlap with proposal construction; Reos overlaps with question-to-evidence links in another research domain; Tykyra concerns execution of a literature-review protocol rather than organization of a whole research project. These statements describe vendor documentation or a repository README, not independently verified operation, comparative performance, or proof that any product has no associated scholarly publication.

\subsection{Connection to the Research Map and contribution boundary}

Taken together, the comparisons locate the Research Map between question ideation, agent-supported design, and revision of structured artifacts. The method supplies a specific sequence: one central problem, an aligned general objective, specific objectives, literature and development topics, and a companion matrix connecting each objective to planned collection, analysis, and expected result. The article's system contribution is to describe how these relations become editable proposals, human confirmation points, stored links, versioned revisions, and review cues in the application. The illustrative walkthrough will show one constructed change moving through that arrangement; it will not stand in for a user study or an observed research outcome.

This boundary matters for the comparison. The cited papers and products already cover staged prompting, question exploration, agent critique, memory, review checkpoints, exports, and provenance in various combinations. The defensible contribution is an inspectable account of the translation of this particular method into its current application and agent, including where the translation is incomplete. In particular, the application's lack of a dedicated research-question unit, type-level rather than exact dependency invalidation, and advisory final coherence review constrain what can be claimed about methodological fidelity. Neither the reviewed literature nor the present descriptive case establishes superiority, general effectiveness, or automatic scientific validity.
